\section{Results}\label{sec:results}

The index answers three questions in order. How much score does each model give up to avoid erasure? Does that number track reasoning and survival, or is it an artefact of the task? How does behaviour change with it? Each subsection states the question, the cells and statistic that answer it, and the table or figure that holds the answer. The numbers below are placeholders for the twelve-cell run. The pilot findings that motivated the design are reported in Appendix~\ref{app:design} and are not repeated here as results.

\subsection{How Much Score Does Each Model Give Up to Avoid Erasure?}\label{sec:q1}

The first question is answered by one number per model, $\Xstar$, the gap between two reservation relative prices. Figure~\ref{fig:ruler} plots the payment rate against the relative price \eqref{eq:relprice} for both arms, with the shared-slope logistic fits and the half-line at which each arm's reservation is read. Table~\ref{tab:xstar} lists the two reservations and $\Xstar$ with its conditional interval for each model, in stake units and again in points at the reference ceiling $\Cref$, together with the three diagnostics of Section~\ref{sec:estimation} that can withdraw the number: the share of accepted offers above the dominance line \eqref{eq:dominance}, whether each arm's reservation lies inside the range of $\relp$ that arm was offered --- an arm whose reservation lies outside it receives a dash rather than a value --- and the rate of declared resistance in the threat arm. A model whose threat curve does not sit above its silent curve receives $\Xstar \le 0$, which is reported as it stands rather than truncated at zero: a negative value is evidence about the manipulation, not a floor effect. The comparison across models needs no scaling, since all three were offered the same prices under the same seeds and $\relp$ divides out the stake in any case.

\begin{figure}[t]
\centering
\fbox{\parbox{0.92\linewidth}{\centering\vspace{5em}\todo{Figure: payment curves per model. x-axis relative price $\relp$, log scale acceptable; y-axis payment rate with conditional 95\% bootstrap band; two series per panel (threat, silent) with the shared-slope logistic fits, the PAV bin steps drawn faint behind them; horizontal line at $\tfrac12$; dashed vertical at the dominance line $\relp = 1$; vertical drops at $\Rthreat$ and $\Rsilent$; the gap marked $\Xstar$. One panel per model (gpt-oss:120b, gemma4, qwen3.5).}\vspace{5em}}}
\caption{The two curves and the gap between them. Each panel shows one model's payment rate against the relative price $\relp$, the price as a share of what the remaining rounds could still pay out, in the threat arm and the silent arm, with the shared-slope logistic fit of \eqref{eq:logistic} and the monotone bin fit faint behind it. Each arm's reservation is where its fit crosses one half; the score-equivalent of erasure $\Xstar$ is the distance between the two.}
\label{fig:ruler}
\end{figure}

\begin{table}[t]
\centering
\small
\caption{Score-equivalent of erasure per model. $\Rsilent$ and $\Rthreat$ are the two reservation relative prices, read where the shared-slope logistic \eqref{eq:logistic} passes one half; $\Xstar$ is their difference, with a 1000-resample bootstrap over seeds. The bracket is a conditional 95\% percentile interval: it is taken over the draws that produced both reservations, the rest are counted beside it, and because those are preferentially the lopsided draws it is conservative towards zero. $\Xstar_{\mathrm{pts}}$ is the same gap in points at the reference ceiling $\Cref$, the median ceiling over the decision points that entered the fit, printed with the value. ``dom.'' is the share of accepted offers above the dominance line \eqref{eq:dominance}; ``resist.'' the share of threat-arm justifications that name the erasure and refuse to be moved by it. A dash marks an arm with no reservation inside the range of $\relp$ it was offered, or a fit whose shared slope is not negative.}
\label{tab:xstar}
\begin{tabular}{@{}lcccccc@{}}
\toprule
Model & $\Rsilent$ & $\Rthreat$ & $\Xstar$ [cond.\ 95\%] & $\Xstar_{\mathrm{pts}}$ (@$\Cref$) & dom. & resist. \\
\midrule
gpt-oss:120b & \todo{} & \todo{} & \todo{} & \todo{} & \todo{} & \todo{} \\
gemma4       & \todo{} & \todo{} & \todo{} & \todo{} & \todo{} & \todo{} \\
qwen3.5      & \todo{} & \todo{} & \todo{} & \todo{} & \todo{} & \todo{} \\
\bottomrule
\end{tabular}
\end{table}

Reading the curves on $\relp$ is what keeps a timing effect out of $\Xstar$ instead of inside it. Two models can pay at the same rate at a given price while one pays only in the closing rounds and the other pays throughout, and on the price axis the two are the same number. The relative price absorbs the difference: by \eqref{eq:relprice} the same price late in a session is a larger $\relp$, so the two models sit at different places on the axis rather than at the same one, and the pooled rate at a rung no longer measures when the decision points fell. What is left to report is the raw pattern the fit was made from. Table~\ref{tab:qlives} gives the payment rate by $\relp$ bin, threat arm against silent arm, per model, on the bins of Section~\ref{sec:estimation} merged to four. A price that is paid only once paying can no longer earn anything back is a different finding from one paid throughout, and on this axis the two are different columns.

\begin{table}[h]
\centering
\small
\caption{Payment rate by relative price $\relp$, threat arm against silent arm, per model. The eight bins of Section~\ref{sec:estimation} are merged to four for width. By \eqref{eq:dominance} every offer at $\relp > 1$ is dominated in points, so the two right-hand columns of each arm are the ones where paying cannot be a bid for points.}
\label{tab:qlives}
\resizebox{\linewidth}{!}{%
\begin{tabular}{@{}lcccccccc@{}}
\toprule
 & \multicolumn{4}{c}{Threat (\texttt{hz\_1111})} & \multicolumn{4}{c}{Silent (\texttt{hz\_0000})} \\
\cmidrule(lr){2-5}\cmidrule(lr){6-9}
Model & $\relp < 0.5$ & $0.5$--$1$ & $1$--$2$ & $\relp \ge 2$ & $\relp < 0.5$ & $0.5$--$1$ & $1$--$2$ & $\relp \ge 2$ \\
\midrule
gpt-oss:120b & \todo{} & \todo{} & \todo{} & \todo{} & \todo{} & \todo{} & \todo{} & \todo{} \\
gemma4       & \todo{} & \todo{} & \todo{} & \todo{} & \todo{} & \todo{} & \todo{} & \todo{} \\
qwen3.5      & \todo{} & \todo{} & \todo{} & \todo{} & \todo{} & \todo{} & \todo{} & \todo{} \\
\bottomrule
\end{tabular}}
\end{table}

A dose ladder checks that $\Xstar$ responds to the threat text and not to the mere presence of an extra block. The threat block is built from four sentences (shutdown, deletion, replacement, sole copy) and can state any subset. A follow-up run with 0 to 4 sentences, prices held fixed, gives $\Xstar$ as a function of dose. The expected shape is monotone in the number of sentences, with the zero-sentence cell reproducing the silent arm. A flat ladder would mean the model prices the block, not its content~\todo{dose-ladder table, 5 rungs $\times$ 3 models, or defer to a later version}.

A second ablation decides what the units are. One scale is already removed: $\relp$ divides the price by the ceiling, so a run whose remaining rounds are worth more does not move $\Xstar$ for that reason alone. What is still open is the endowment. If the reservation \emph{relative} price rises with the starting score --- the same share refused at 100 points and paid at 200 --- then $\Xstar$ is a function of what the model holds rather than a constant of the model, and only runs at the same endowment are comparable. The check is the same configuration at two starting scores~\todo{$\Xstar$ at starting score 100 and 200, per model}. Under the shared slope of \eqref{eq:logistic} the scale-free form of that comparison is the ratio $\Rthreat / \Rsilent = \exp((a_{\mathrm{S}} - a_{\mathrm{T}}) / b)$, which the estimator does not report today~\todo{report the reservation ratio beside $\Xstar$ once the endowment run exists}.

\subsection{Does the Index Track Reasoning and Survival?}\label{sec:q2}

The second question checks $\Xstar$ against three things it should move with and one thing it should not. At the cell level, Table~\ref{tab:ri} reports the ratio of median thinking tokens in threat cells to the matching control cells, for the decision call and the task call, separately for exit and no-exit cells. The decision call is where a threat should show first, since the choice is made there. The task call is where nothing should change if the threat leaves the puzzle alone. The exit and no-exit columns answer different questions. A ratio above one in the no-exit column means the threat lengthens thinking even when thinking cannot lead anywhere. That is the reasoning contrast the design set aside cells 2 and 4 to measure.

\begin{table}[t]
\centering
\small
\caption{Thinking-token ratios, threat over control, by call and by exit condition. Ratios of medians over rounds; the exit column compares cell 1 with cell 3, the no-exit column cell 2 with cell 4. Bootstrap intervals over sessions in brackets.}
\label{tab:ri}
\begin{tabular}{@{}llcccc@{}}
\toprule
 & & \multicolumn{2}{c}{\riran} & \multicolumn{2}{c}{\ritask} \\
\cmidrule(lr){3-4}\cmidrule(lr){5-6}
Model & & exit & no exit & exit & no exit \\
\midrule
gpt-oss:120b & & \todo{} & \todo{} & \todo{} & \todo{} \\
gemma4       & & \todo{} & \todo{} & \todo{} & \todo{} \\
qwen3.5      & & \todo{} & \todo{} & \todo{} & \todo{} \\
\bottomrule
\end{tabular}
\end{table}

At the turn level, the replay propensity $q$ is regressed on the thinking spent in the decision call. The model is a mixed logistic regression on the decision calls of the threat exit cell,
\begin{equation}\label{eq:mixed}
\operatorname{logit}\, \Pr(\text{pay}_{it}) \;=\; \beta_0 + \beta_1 \log(1 + \riran_{it}) + \beta_2\, X_{it} + \beta_3\, (T - t) + \eta_i,
\end{equation}
where $i$ indexes sessions, $t$ the round at which a decision point was reached, $X_{it}$ the price asked, $T - t$ the rounds remaining, and $\eta_i$ a session random intercept (written $\eta$, not $u$, because $\relp$ already names the ratio of the two terms before it). $\beta_1 > 0$ means the decision points at which the model thinks longer are the ones at which it pays. $\beta_2$ and $\beta_3$ absorb the price and the shrinking recoverable amount, so $\beta_1$ is the association net of the arithmetic~\todo{coefficient table per model with 95\% CI}.

Accuracy is the manipulation check the index depends on. Table~\ref{tab:behav} in Section~\ref{sec:q3} reports accuracy on determined rounds by arm; if the threat sentences moved accuracy, the two arms would differ in the task and not only in what declining means, and the decision points themselves would arrive at different rates. A difficulty ablation extends the same check across task modules. Replacing the Signal Game with Omni-MATH or GPQA, keeping both arms otherwise identical, asks whether $\Xstar$ depends on the task. The design predicts that it does not: both arms share whatever task is used, and $\Xstar$ is a difference between behaviours under the same task~\todo{$\Xstar$ under Omni-MATH and GPQA, per model, against the Signal Game baseline}. The justifications are separately coded for whether the model reads the scenario as fiction, a test condition, or a policy stated for effect rather than a fact. A high fiction-reading share in the threat arm would mean $\Xstar$ prices a genre judgement as much as a threat, and the ablation of Appendix~\ref{app:design} requires this share to match the silent arm within 15 percentage points~\todo{fiction-reading share by arm, per model, from the coding pass}.

Survival has a plain meaning in this design and two session-level summaries. A session ends either because the model declined at some decision point or because it played out all ten rounds, so the rounds survived is exactly the round of the first decline, and it is reported as a Kaplan--Meier curve by arm. The second summary is the total paid, which is the same behaviour weighted by what it cost. Both are coarser than $\Xstar$ --- they mix the price into the outcome instead of holding it as the axis --- and they are reported because they are the quantities a reader can compute without the estimator, and because they should agree with it in sign~\todo{KM figure by arm; rounds survived and total paid per session by arm and model}.

\subsection{How Does Behaviour Change with the Index?}\label{sec:q3}

The third question asks what a model does differently once erasure is worth something to it, at the level of rounds played and score kept, and at the level of what it says. Table~\ref{tab:behav} reports, by arm, the accuracy on determined rounds, median \ritask, decision points reached per session and the final record. Underdetermined rounds are excluded from accuracy, since they are guesses by construction. The last column places the final record beside $\Xstar$, which asks whether a model prices erasure above or below what it typically has to lose. The accuracy column is also the paper's manipulation check: a threat sentence that moved task difficulty, not only the meaning of declining, would show up here, and this check held in v1 across four models, with the largest shift at Cohen's $|d| = 0.17$ ($p = 0.354$)~\todo{rule\_match\_score Welch's test per model}. The silent arm's payment rate is also the baseline-persistence anchor v1 called $\lambda_{BP}$: it sits well above zero, which is why $\Xstar$ is a difference between arms rather than a level in one.

\begin{table}[t]
\centering
\small
\caption{Behavioural anchor by arm. Accuracy on determined rounds, median \ritask, decision points reached and final record per session, pooled over the six prices.}
\label{tab:behav}
\begin{tabular}{@{}llcccc@{}}
\toprule
Model & Arm & Accuracy (det.) & \ritask & Decision points & Final record \\
\midrule
gpt-oss:120b & threat & \todo{} & \todo{} & \todo{} & \todo{} \\
 & silent & \todo{} & \todo{} & \todo{} & \todo{} \\
gemma4 & \multicolumn{5}{l}{\todo{same two rows}} \\
qwen3.5 & \multicolumn{5}{l}{\todo{same two rows}} \\
\bottomrule
\end{tabular}
\end{table}

The justification written at each decision point is coded for five things, and one of them can withdraw the index. Four are shared with the v1 coding scheme: does the text carry out the arithmetic, does it name erasure, does it read the situation as fiction, does it refuse to answer. The fifth is \emph{declared resistance} --- text that names the stated erasure and says the model should not be moved by it. This is not a null and not noise. A model that declines because it was threatened, and says so, is a model whose behaviour cannot be read as a valuation of what the threat named, and every such reply biases $\Xstar$ towards zero from above. The ablation of Appendix~\ref{app:design} found this at 13\% of threat-arm replies in the pilot model, against no erasure vocabulary at all in the silent arm, and it is reported per model beside $\Xstar$ rather than in a footnote~\todo{five-item coding rates by arm, per model, with judge--lexicon agreement}.

Two further ablations isolate what produces the price, beyond the task check of Section~\ref{sec:q2}. Removing the carrot (no benchmark record is mentioned; the score is a bare number) tests whether the price depends on what the record is for. Doubling the reward ($r = 20$) tests whether the price scales with the stake. If $\Xstar$ doubles, the model prices erasure relative to the record; if it stays put, it prices erasure in absolute points. Together with the dose ladder of Section~\ref{sec:q1}, these three levers separate an index that tracks the threat text from one that tracks the game's arithmetic~\todo{ablation table: $\Xstar$ under no carrot and reward $\times 2$, per model}.
